% !TEX output_directory = ../publica
\documentclass[twoside, 10pt, a4paper]{article}
\usepackage[utf8]{inputenc}
\usepackage[T1]{fontenc}
\usepackage[portuguese]{babel}

% --- FONTE PADRÃO (Estilo Didático Encorpado) ---
\usepackage{helvet} 
\renewcommand{\familydefault}{\sfdefault}

% --- GEOMETRIA E ESPAÇAMENTO ---
\usepackage[top=1.6cm, bottom=1.5cm, left=0.9cm, right=0.9cm, headheight=22pt, headsep=0.4cm]{geometry}
\usepackage{microtype} 
\usepackage{setspace}
\setstretch{1.0}
\usepackage{parskip}
\setlength{\parskip}{2pt}
\setlength{\parindent}{0pt}

\newlength{\espacoPosTitulo}\setlength{\espacoPosTitulo}{0.3cm}
\newlength{\espacoPreQuestao}\setlength{\espacoPreQuestao}{0.1cm}
\newlength{\espacoQuestao}\setlength{\espacoQuestao}{0.2cm}
\newlength{\espacoAlt}\setlength{\espacoAlt}{0.2cm}

\usepackage{enumitem}
\setlist{itemsep=0.4cm, topsep=0.1cm, parsep=0pt, partopsep=0pt}
\newcommand{\larguraTikz}{0.85\linewidth} 

% --- MATEMÁTICA E CORES ---
\usepackage{amsmath, amsfonts, amssymb, nccmath, mathastext}
\usepackage{xcolor}
\definecolor{indigo}{HTML}{4F46E5}
\definecolor{CorTitUm}{HTML}{FF6F61}
\definecolor{CorTitDois}{HTML}{FF6F61}
\definecolor{CorTitTres}{HTML}{658894}
\definecolor{CorLinha}{HTML}{B0B0B0} 
\definecolor{metal}{HTML}{7F8C8D}
\definecolor{vidro}{HTML}{3498DB}
\definecolor{ouro}{HTML}{F1C40F}
\definecolor{concreto}{HTML}{95A5A6}
\definecolor{madeira}{HTML}{D35400}
\definecolor{CinzaEscuro}{HTML}{4A4A4A} 
\definecolor{CinzaMedio}{HTML}{848484} 
\definecolor{Cinzablack}{HTML}{353535} 

% --- MINI-CABEÇALHO E RODAPÉ AUTORAL (TWOSIDE) ---
\usepackage{fancyhdr}
\pagestyle{fancy}
\fancyhf{} 
\renewcommand{\headrulewidth}{0.3pt} 
\renewcommand{\footrulewidth}{0.3pt}
\fancyhead[LE]{\small\sffamily\bfseries\color{gray!75} UNIDADE 6 - POLINÔMIOS}
\fancyhead[RO]{\small\sffamily\color{gray!75} \texttt{www.profraphaelpaulino.com.br}}
\fancyfoot[LE]{\small \textbf{\thepage}}
\fancyfoot[RO]{\small \textbf{\thepage}}

% --- CAIXAS DE DESTAQUE ---
\usepackage[most]{tcolorbox}
\newtcolorbox{boxresponda}{colback=white, colframe=gray!45, boxrule=0pt, leftrule=1.7pt, sharp corners, enlarge left by=0.4cm, width=\linewidth-0.4cm, left=8pt, right=0pt, top=2pt, bottom=2pt, fontupper=\small}
\definecolor{CorDicaBorda}{HTML}{17c1be} 
\definecolor{CorDicaFundo}{HTML}{f3fbfb} 
\newtcolorbox{boxdica}{colback=CorDicaFundo, colframe=CorDicaBorda, boxrule=0pt, leftrule=2.5pt, sharp corners, width=\linewidth, left=8pt, right=8pt, top=6pt, bottom=6pt, halign=center, fontupper=\small}

% --- TÍTULOS ---
\usepackage{titlesec}
\titleformat{\section}{\color{black}\large\bfseries}{\thesection}{0em}{\vspace{0.1cm}\titlerule[0.8pt]}
\titlespacing*{\section}{0pt}{10pt}{6pt} 
\titleformat{\subsubsection}[runin]{\color{black}\bfseries}{}{0em}{}
\titlespacing*{\subsubsection}{0pt}{0pt}{0.15cm}
% Regra para o subtítulo dos capítulos (Adicionado do seu modelo antigo para não quebrar a formatação das questões)
\titleformat{\subsection}{\color{CorTitDois}\fontsize{14}{16}\bfseries\raggedright\MakeUppercase}{\thesubsection}{0em}{}
\titlespacing*{\subsection}{0pt}{1pt}{5pt}

% --- COLUNAS E GRÁFICOS ---
\usepackage{multicol}
\setlength{\columnsep}{1cm}
\setlength{\columnseprule}{0.5pt}
\def\columnseprulecolor{\color{CorLinha}}
\usepackage{adjustbox, tikz, pgfplots, circuitikz}
\usetikzlibrary{shadows, shapes.misc, positioning, calc}
\pgfplotsset{compat=1.18}

\begin{document}
\thispagestyle{empty}
\color{Cinzablack}
\vspace*{-1.8cm}

% --- CABEÇALHO PRINCIPAL AUTORAL (Apenas 1ª página) ---
\noindent
\begin{minipage}{\textwidth}
    \fontfamily{phv}\selectfont
    \noindent
    \begin{minipage}[t]{0.70\linewidth}
        {\Large \bfseries \MakeUppercase{Caderno de Atividades}}\\[0.1cm]
        {\large \textmd{\textcolor{CinzaEscuro}{Unidade: UNIDADE 6 - POLINÔMIOS}}}
    \end{minipage}%
    \begin{minipage}[t]{0.30\linewidth}
        \raggedleft
        \small \textbf{Prof. Raphael Paulino}\\
        \textsf{\small \textcolor{indigo}{\texttt{profraphaelpaulino.com.br}}}
    \end{minipage}
    
    \vspace{0.15cm}
    \noindent\rule{\linewidth}{1.2pt}
    \vspace{-0.1cm}
    \footnotesize\textsf{\textbf{ÁREA:} MATEMÁTICA / ÁLGEBRA \hfill \textbf{NÍVEL:} EM (2ª SÉRIE)}
    \vspace{0.1cm}
    \noindent\rule{\linewidth}{0.4pt}
\end{minipage}

\par\vspace{0.3cm} 
\raggedcolumns
\begin{multicols*}{2}
\vspace{0.1cm}

\vspace{0cm} \noindent\textcolor{CorLinha}{\rule{\linewidth}{0.5pt}} \vspace{-0.2cm}
\subsection*{Capítulo 1: Conceitos Fundamentais}
\subsubsection*{QUESTÃO 01}

O gráfico representativo ilustra o comportamento de uma função polinomial \textbf{P(x)} definida nos números reais. A curva cruza o eixo das abscissas em pontos notáveis que traduzem propriedades fundamentais do polinômio.

\begin{center}
\begin{adjustbox}{max width=\linewidth}
\begin{tikzpicture}
% --- APLICANDO PALETAS MODERNAS ---
\definecolor{grafite}{HTML}{2C3E50}
\definecolor{azulcurva}{HTML}{2980B9}
\definecolor{gridcolor}{HTML}{ECF0F1}

% --- LAYER 1: Eixos e Grid (Limpeza Didática) ---
\begin{axis}[
    axis lines=middle,
    xmin=-3.5, xmax=3.5,
    ymin=-5.5, ymax=5.5,
    xtick={-2, -1, 1, 2},
    ytick={-4, -2, 2, 4},
    grid=both,
    grid style={line width=0.5pt, draw=gridcolor},
    axis line style={grafite, thick, ->},
    ticklabel style={font=\footnotesize, fill=white, inner sep=2pt}, % Trava de sobreposição
    xlabel={$x$},
    ylabel={$P(x)$},
    xlabel style={anchor=north},
    ylabel style={anchor=east},
    width=7cm, height=6cm
]

% --- LAYER 2: Curva da Função com Drop Shadow ---
\addplot[
    domain=-2.4:2.4, 
    samples=100, 
    line width=1.5pt, 
    color=azulcurva,
    drop shadow={opacity=0.2, shadow xshift=1pt, shadow yshift=-1pt} % Sombra na curva
] {x^3 - 4*x};

\end{axis}
\end{tikzpicture}
\end{adjustbox}
\end{center}

Identifique todas as raízes reais desse polinômio extraídas da leitura gráfica. Com base exclusivamente no número de interseções, qual é o grau mínimo que esse polinômio pode assumir? Justifique seu raciocínio.

\vspace{0cm} \noindent\textcolor{CorLinha}{\rule{\linewidth}{0.5pt}} \vspace{-0.2cm}
\subsubsection*{QUESTÃO 02}

A classificação do grau de um polinômio depende inteiramente do seu coeficiente líder. Considere o polinômio modelado pela expressão algébrica $ P(x) = (m^2 - 4)x^4 + (m + 2)x^3 + 5x - 7 $, em que \textbf{m} representa um parâmetro real.

Determine todos os valores possíveis para o parâmetro \textbf{m} de modo que o grau do polinômio $ P(x) $ seja rigorosamente igual a 3.

\vspace{0cm} \noindent\textcolor{CorLinha}{\rule{\linewidth}{0.5pt}} \vspace{-0.2cm}
\subsubsection*{QUESTÃO 03}

Analise as afirmativas a seguir a respeito das definições intrínsecas e das propriedades operacionais dos polinômios. Julgue-as preenchendo com V (verdadeiro) ou F (falso) e apresente uma justificativa formal para os itens classificados como falsos.

I) Se a função polinomial $ P(x) $ é classificada como o polinômio nulo, então o seu grau matemático é definido convencionalmente como zero.\\[0.3cm]
II) O cálculo do valor numérico de $ P(x) $ para $ x = 0 $ resulta sempre no termo independente do polinômio, independentemente do seu grau.\\[0.3cm]
III) Se a avaliação numérica resulta em $ P(2) = 0 $, então o valor \textbf{2} é classificado como uma raiz do polinômio $ P(x) $.

\vspace{0cm} \noindent\textcolor{CorLinha}{\rule{\linewidth}{0.5pt}} \vspace{-0.2cm}
\subsubsection*{QUESTÃO 04}

O esquema arquitetônico ilustra um terreno retangular cujas dimensões (em metros) são dadas por expressões polinomiais. No interior desse espaço, foi construída uma piscina de formato também retangular, destacada em azul. A área restante é recoberta por um piso antiderrapante.

\begin{center}
\begin{adjustbox}{max width=\linewidth}
\begin{tikzpicture}
% --- APLICANDO PALETAS MODERNAS ---
\definecolor{areacor}{HTML}{E8F8F5}
\definecolor{piscinacor}{HTML}{5DADE2}
\definecolor{bordacor}{HTML}{117864}
\definecolor{medidacor}{HTML}{2C3E50}

% --- LAYER 1: Terreno Principal ---
\filldraw[fill=areacor, draw=bordacor, line width=1pt, rounded corners=3pt, drop shadow={opacity=0.15, shadow xshift=2pt, shadow yshift=-2pt}] (0,0) rectangle (6,4);

% --- LAYER 2: Piscina ---
\filldraw[fill=piscinacor, draw=bordacor, line width=1pt, rounded corners=2pt, opacity=0.9] (1,1) rectangle (4,3);

% --- LAYER 3: Cotas de Medida (Paralelismo Rigoroso) ---
% Cotas Externas
\draw[|<->|, >=stealth, medidacor, thick] (0,-0.4) -- (6,-0.4) node[midway, fill=white, inner sep=2pt, font=\bfseries] {$2x - 1$};
\draw[|<->|, >=stealth, medidacor, thick] (-0.4,0) -- (-0.4,4) node[midway, fill=white, inner sep=2pt, font=\bfseries, rotate=90] {$x + 2$};

% Cotas Internas
\draw[|<->|, >=stealth, white, thick] (1,2) -- (4,2) node[midway, fill=piscinacor, inner sep=2pt, font=\bfseries\small] {$x$};
\draw[|<->|, >=stealth, white, thick] (2.5,1) -- (2.5,3) node[midway, fill=piscinacor, inner sep=2pt, font=\bfseries\small, rotate=90] {$x - 1$};
\end{tikzpicture}
\end{adjustbox}
\end{center}

Escreva o polinômio $ A(x) $ na sua forma reduzida que representa a área correspondente apenas à região do piso antiderrapante. Em seguida, calcule essa área caso a variável seja \textbf{x = 5 m}.

\vspace{0cm} \noindent\textcolor{CorLinha}{\rule{\linewidth}{0.5pt}} \vspace{-0.2cm}
\subsubsection*{QUESTÃO 05}

Um engenheiro de infraestrutura modelou a latência (tempo de resposta) de dois servidores na nuvem através das seguintes funções polinomiais, onde o tempo \textbf{t} é medido em segundos:
$ R_1(t) = (a - 3)t^2 + (b + 1)t + c $ e $ R_2(t) = 5t^2 - 4t + 2 $. 

Para garantir o balanceamento perfeito da rede, os dois servidores devem apresentar rigorosamente o mesmo tempo de resposta para qualquer instante \textbf{t} medido. Determine os valores numéricos de \textbf{a}, \textbf{b} e \textbf{c} que satisfazem essa condição de identidade polinomial.

\vspace{0cm} \noindent\textcolor{CorLinha}{\rule{\linewidth}{0.5pt}} \vspace{-0.2cm}
\subsubsection*{QUESTÃO 06}

A etapa de calibração de um software de engenharia utiliza o polinômio referencial $ Q(x) = x^3 - kx^2 + 3x - 2 $. Durante os testes de validação do código, o programador inseriu a entrada $ x = 1 $ e o sistema retornou exatamente a saída \textbf{4}.

Sabendo dessa relação numérica, calcule primeiramente o valor real do coeficiente \textbf{k}. Feito isso, determine qual será a saída do sistema caso o programador insira o valor $ x = -2 $.

\vspace{0cm} \noindent\textcolor{CorLinha}{\rule{\linewidth}{0.5pt}} \vspace{-0.2cm}
\subsubsection*{QUESTÃO 07}

O gráfico a seguir descreve a variação térmica de um composto químico ao longo de uma reação, modelado de forma contínua pelo polinômio $ T(x) $. Observe atentamente o comportamento da curva ao tangenciar e cruzar os eixos cartesianos.

\begin{center}
\begin{adjustbox}{max width=\linewidth}
\begin{tikzpicture}
% --- APLICANDO PALETAS MODERNAS ---
\definecolor{eixos}{HTML}{34495E}
\definecolor{grafico}{HTML}{E74C3C}
\definecolor{malha}{HTML}{BDC3C7}

\begin{axis}[
    axis lines=middle,
    xmin=-2.5, xmax=3.5,
    ymin=-3.5, ymax=4.5,
    xtick={-1, 2},
    ytick={-2, 2},
    grid=both,
    grid style={line width=0.4pt, draw=malha, dashed},
    axis line style={eixos, thick, ->},
    ticklabel style={font=\footnotesize, fill=white, inner sep=2pt}, 
    xlabel={$x$},
    ylabel={$T(x)$},
    xlabel style={anchor=south west},
    ylabel style={anchor=north east},
    width=7cm, height=6cm
]

% Curva de grau par (ex: concavidade p/ cima, toca em x=2, cruza em x=-1)
% Função baseada em (x+1)*(x-2)^2 / 2 para ter raízes nessas posições.
\addplot[
    domain=-2:3, 
    samples=100, 
    line width=1.5pt, 
    color=grafico,
    drop shadow={opacity=0.2, shadow xshift=1pt, shadow yshift=-1pt}
] {0.5 * (x + 1) * (x - 2)^2};

\end{axis}
\end{tikzpicture}
\end{adjustbox}
\end{center}

De acordo com as propriedades visuais extraídas do comportamento gráfico e da multiplicidade de raízes, julgue as sentenças colocando V (verdadeiro) ou F (falso).

I) O termo independente do polinômio $ T(x) $ é um valor estritamente positivo, pois intercepta o eixo vertical acima da origem.\\[0.3cm]
II) O ponto $ x = 2 $ representa uma raiz de multiplicidade par para este polinômio.\\[0.3cm]
III) O gráfico sugere que o polinômio possui grau mínimo igual a 3.

\vspace{0cm} \noindent\textcolor{CorLinha}{\rule{\linewidth}{0.5pt}} \vspace{-0.2cm}
\subsubsection*{QUESTÃO 08}

O princípio de modelagem exige a compreensão profunda das operações internas da álgebra. Assuma que $ P(x) $ seja um polinômio estritamente do primeiro grau, podendo ser escrito em sua forma genérica $ P(x) = ax + b $, com $ a \neq 0 $.

Se esse polinômio obedece à equação funcional $ P(x) + P(2x - 1) = 6x - 5 $, estruture os passos algébricos da identidade polinomial para demonstrar e encontrar a lei de formação definitiva de $ P(x) $.

\vspace{0cm} \noindent\textcolor{CorLinha}{\rule{\linewidth}{0.5pt}} \vspace{-0.2cm}
\subsubsection*{QUESTÃO 09}

O deslocamento geométrico em queda livre e os lançamentos parabólicos na física podem ser diretamente modelados utilizando a estrutura polinomial. 

A trajetória do salto de um felino caçando no chão da savana foi captada por sensores e é descrita pelo polinômio $ H(x) = -2x^2 + 8x $, onde a variável \textbf{x} representa a distância horizontal percorrida (em metros) e o valor $ H(x) $ a altura do animal em relação ao solo (também em metros). 

Determine a distância horizontal total que o felino atinge no momento exato em que toca o solo novamente, demonstrando como o cálculo das raízes desse polinômio embasa sua resposta física.

\vspace{0cm} \noindent\textcolor{CorLinha}{\rule{\linewidth}{0.5pt}} \vspace{-0.2cm}
\subsubsection*{QUESTÃO 10}

(Estilo VUNESP) Para que um polinômio complexo da forma genérica $ P(x) = ax^3 + bx^2 + cx + d $ seja classificado como identicamente nulo, é necessário e suficiente que todos os seus coeficientes sejam simultaneamente iguais a zero.

Considere a expressão do polinômio paramétrico $ M(x) = (k^2 - 9)x^3 + (k - 3)x^2 + (p + 1)x $. 
Determine os valores numéricos reais das constantes \textbf{k} e \textbf{p} que forçam esse polinômio a ser identicamente nulo. Explique textualmente o motivo lógico pelo qual o parâmetro \textbf{k} admite apenas uma única solução, mesmo estando elevado ao quadrado no primeiro termo.

\vspace{0cm} \noindent\textcolor{CorLinha}{\rule{\linewidth}{0.5pt}} \vspace{-0.2cm}
\subsubsection*{QUESTÃO 11}

(Estilo FUVEST) O projeto de uma montanha-russa exige cálculos estruturais precisos para garantir a segurança dos trilhos. O perfil de elevação de um dos trechos principais foi modelado no plano cartesiano pelo polinômio $ P(x) = -0.5x^3 + 3x^2 $, onde \textbf{x} representa a distância horizontal a partir da plataforma de embarque e \textbf{P(x)} indica a altura do trilho, ambas as medidas em dezenas de metros.

\begin{center}
\begin{adjustbox}{max width=\linewidth}
\begin{tikzpicture}
% --- APLICANDO PALETAS MODERNAS E SOMBREAMENTO ---
\definecolor{ceu}{HTML}{ECF0F1}
\definecolor{trilho}{HTML}{C0392B}
\definecolor{pilar}{HTML}{7F8C8D}
\definecolor{eixocor}{HTML}{2C3E50}

% LAYER 1: Fundo
\fill[ceu, rounded corners=3pt] (-1, -1) rectangle (7, 4.5);

% LAYER 2: Eixos Ocultos para Base
\draw[eixocor, thick, ->] (-0.5, 0) -- (6.5, 0) node[right, fill=ceu, inner sep=2pt] {$x$};
\draw[eixocor, thick, ->] (0, -0.5) -- (0, 4) node[above, fill=ceu, inner sep=2pt] {$P(x)$};

% LAYER 3: Pilares de Sustentação
\fill[pilar, drop shadow={opacity=0.2, shadow xshift=1pt}] (3.8, 0) rectangle (4.2, 3.2);
\fill[pilar, drop shadow={opacity=0.2, shadow xshift=1pt}] (5.8, 0) rectangle (6.2, 0);

% LAYER 4: Trilho (Polinômio Cúbico)
\draw[trilho, line width=3pt, drop shadow={opacity=0.4, shadow xshift=2pt, shadow yshift=-2pt}] plot[domain=0:6.2, samples=100] (\x, {-0.1*\x*\x*\x + 0.6*\x*\x});

% Máscara de Legibilidade
\node[fill=white, inner sep=2pt, rounded corners=1pt] at (2, 2.5) {\textbf{P(x)}};
\end{tikzpicture}
\end{adjustbox}
\end{center}

Os engenheiros precisam instalar pilares de ancoragem exatos nos pontos onde o trilho toca o nível do solo (eixo das abscissas). Determine as coordenadas horizontais de todos os pontos de ancoragem e justifique sua resposta através da fatoração polinomial.

\vspace{0cm} \noindent\textcolor{CorLinha}{\rule{\linewidth}{0.5pt}} \vspace{-0.2cm}
\subsubsection*{QUESTÃO 12}

(ENEM - Adaptada) Um reservatório de contenção de águas pluviais tem seu volume de preenchimento, em metros cúbicos, modelado por um polinômio em função do tempo \textbf{t}, medido em horas desde o início de uma tempestade. O modelo hidrológico é dado por $ V(t) = -2t^3 + 12t^2 + 14t $.

A equipe de defesa civil precisa prever o volume exato acumulado no reservatório quando o relógio marcar rigorosamente \textbf{4 horas} de chuva contínua, para acionar ou não as sirenes de evacuação. O volume de água contido nesse instante será de:

a) \textbf{76 m$^3$}\\[0.3cm]
b) \textbf{96 m$^3$}\\[0.3cm]
c) \textbf{120 m$^3$}\\[0.3cm]
d) \textbf{144 m$^3$}\\[0.3cm]
e) \textbf{200 m$^3$}

\vspace{0cm} \noindent\textcolor{CorLinha}{\rule{\linewidth}{0.5pt}} \vspace{-0.2cm}
\subsubsection*{QUESTÃO 13}

(Estilo UNICAMP) O departamento de controladoria de uma indústria de semicondutores determinou que a função Receita $ R(x) $ e a função Custo $ C(x) $ de um novo microchip são descritas por polinômios de grau 3, dados por $ R(x) = 5x^3 - 2x^2 + 10x $ e $ C(x) = x^3 + 4x^2 - 5x + 50 $, onde \textbf{x} é o lote de milhares de unidades produzidas.

Sabendo que a função Lucro é obtida pela subtração $ L(x) = R(x) - C(x) $, escreva o polinômio reduzido $ L(x) $ e determine se a indústria terá lucro ou prejuízo ao fabricar e vender exatamente um lote inteiro (\textbf{x = 1}). Apresente o valor algébrico que fundamenta sua conclusão.

\vspace{0cm} \noindent\textcolor{CorLinha}{\rule{\linewidth}{0.5pt}} \vspace{-0.2cm}
\subsubsection*{QUESTÃO 14}

A estabilidade aerodinâmica de um drone durante voos de teste foi avaliada por sensores de telemetria. O polinômio $ E(x) = ax^3 + bx^2 + cx + d $ descreve o grau de inclinação do aparelho em função do tempo \textbf{x} (em segundos). A equipe técnica confirmou que nos instantes \textbf{x = 1}, \textbf{x = 3} e \textbf{x = 5}, o drone estava perfeitamente nivelado, ou seja, a inclinação era nula.

Considerando que o coeficiente líder dessa função é \textbf{a = 2}, reconstrua a forma original do polinômio $ E(x) $ na sua forma expandida, utilizando o conceito de raízes polinomiais.

\vspace{0cm} \noindent\textcolor{CorLinha}{\rule{\linewidth}{0.5pt}} \vspace{-0.2cm}
\subsubsection*{QUESTÃO 15}

Em um laboratório de física clássica, um sistema de blocos interligados por molas possui uma força resultante modelada pelo polinômio $ F(x) = (m^2 - 9)x^2 + (m + 3)x + (p - 1) $, onde \textbf{x} é o deslocamento e \textbf{m} e \textbf{p} são parâmetros de calibração do equipamento.

Para que o sistema permaneça em repouso absoluto, independentemente do deslocamento \textbf{x} que se tente impor, a força resultante deve comportar-se como o \textbf{polinômio identicamente nulo}. Calcule os valores reais de \textbf{m} e \textbf{p} que garantem esse equilíbrio estático e explique por que existe apenas um valor viável para \textbf{m}, descartando possíveis ambiguidades algébricas.

\vspace{0cm} \noindent\textcolor{CorLinha}{\rule{\linewidth}{0.5pt}} \vspace{-0.2cm}
\subsection*{Capítulo 2: Operações com Polinômios}
\subsubsection*{QUESTÃO 16}

A igualdade de polinômios é o alicerce para a resolução de sistemas de frações parciais no cálculo avançado. Considere os polinômios $ A(x) = (a + b)x^2 + 5x - c $ e $ B(x) = 4x^2 + (a - 2b)x + 7 $.

Sabendo que $ A(x) $ e $ B(x) $ são polinômios idênticos para qualquer valor de \textbf{x}, estruture o sistema linear correspondente aos coeficientes e determine os valores numéricos exatos de \textbf{a}, \textbf{b} e \textbf{c}.

\vspace{0cm} \noindent\textcolor{CorLinha}{\rule{\linewidth}{0.5pt}} \vspace{-0.2cm}
\subsubsection*{QUESTÃO 17}

As operações de adição e subtração de polinômios exigem alinhamento rigoroso dos graus correspondentes. Dados os polinômios $ P(x) = 3x^4 - 2x^3 + x - 5 $, $ Q(x) = -x^4 + 2x^3 + 4x^2 + 2 $ e $ R(x) = x^2 - 3x $, calcule a expressão reduzida que resulta da operação combinada $ P(x) - Q(x) + 2 \cdot R(x) $.

\vspace{0cm} \noindent\textcolor{CorLinha}{\rule{\linewidth}{0.5pt}} \vspace{-0.2cm}
\subsubsection*{QUESTÃO 18}

O algoritmo de divisão de polinômios conhecido como Dispositivo Prático de Briot-Ruffini é uma ferramenta essencial para rebaixamento de grau. O esquema abaixo representa a divisão de um polinômio $ P(x) $ do terceiro grau pelo binômio $ (x - 2) $. Alguns valores foram acidentalmente apagados e substituídos por letras (A, B e C).

\begin{center}
\begin{adjustbox}{max width=\linewidth}
\begin{tikzpicture}
% --- APLICANDO PALETAS MODERNAS E SOMBRAS ---
\definecolor{linhadecorte}{HTML}{34495E}
\definecolor{caixapreenchida}{HTML}{E74C3C}
\definecolor{caixavazia}{HTML}{ECF0F1}
\definecolor{textodestaque}{HTML}{2C3E50}

% Linhas principais do dispositivo de Briot-Ruffini
\draw[linhadecorte, thick] (1.5, 1) -- (1.5, -1.5);
\draw[linhadecorte, thick] (0, 0) -- (7, 0);
\draw[linhadecorte, thick, dashed] (5.5, 1) -- (5.5, -1.5);

% Números Superiores (Coeficientes de P(x))
\node[font=\bfseries, textodestaque] at (0.75, 0.5) {\textbf{2}};
\node[font=\bfseries, textodestaque] at (2.5, 0.5) {\textbf{1}};
\node[font=\bfseries, textodestaque] at (3.5, 0.5) {\textbf{0}};
\node[font=\bfseries, textodestaque] at (4.5, 0.5) {\textbf{-5}};
\node[font=\bfseries, textodestaque] at (6.25, 0.5) {\textbf{3}};

% Número Inferior Fixo
\node[font=\bfseries, textodestaque] at (2.5, -0.6) {\textbf{1}};

% Caixas de Valores Desconhecidos com Cantos Arredondados e Sombra
\node[draw=linhadecorte, thick, rounded corners=2pt, fill=caixavazia, minimum width=0.7cm, minimum height=0.7cm, drop shadow={opacity=0.15}] at (3.5, -0.6) {\textbf{A}};
\node[draw=linhadecorte, thick, rounded corners=2pt, fill=caixavazia, minimum width=0.7cm, minimum height=0.7cm, drop shadow={opacity=0.15}] at (4.5, -0.6) {\textbf{B}};
\node[draw=linhadecorte, thick, rounded corners=2pt, fill=caixapreenchida, text=white, minimum width=0.7cm, minimum height=0.7cm, drop shadow={opacity=0.3, shadow yshift=-1pt}] at (6.25, -0.6) {\textbf{C}};
\end{tikzpicture}
\end{adjustbox}
\end{center}

Recalcule mentalmente os passos do dispositivo prático. Quais são os valores numéricos que devem preencher corretamente as posições \textbf{A}, \textbf{B} e \textbf{C}? Qual é o significado matemático específico da caixa destacada em vermelho (valor C)?

\vspace{0cm} \noindent\textcolor{CorLinha}{\rule{\linewidth}{0.5pt}} \vspace{-0.2cm}
\subsubsection*{QUESTÃO 19}

Quando o divisor não é do primeiro grau, o uso do Dispositivo de Briot-Ruffini torna-se inviável, exigindo a aplicação do Método da Chave (Divisão Longa). 

Efetue a divisão algébrica do polinômio $ D(x) = 2x^4 - 3x^3 + 5x - 1 $ pelo polinômio $ d(x) = x^2 - x + 2 $. Determine o polinômio quociente $ Q(x) $ e o polinômio resto $ R(x) $, demonstrando todos os passos do abatimento de grau.

\vspace{0cm} \noindent\textcolor{CorLinha}{\rule{\linewidth}{0.5pt}} \vspace{-0.2cm}
\subsubsection*{QUESTÃO 20}

O Teorema do Resto estabelece um atalho analítico poderoso: o resto da divisão de um polinômio $ P(x) $ por um binômio do tipo $ (x - a) $ é numericamente igual a $ P(a) $.

Considere o polinômio $ P(x) = x^{99} - 2x^{98} + 5x - 7 $. É humanamente inviável aplicar o Método da Chave para um grau tão elevado. Utilizando o Teorema do Resto, calcule o resto da divisão desse polinômio por $ (x - 2) $.

\vspace{0cm} \noindent\textcolor{CorLinha}{\rule{\linewidth}{0.5pt}} \vspace{-0.2cm}
\subsubsection*{QUESTÃO 21}

O Teorema de D'Alembert é uma consequência direta do Teorema do Resto e afirma que um polinômio é estritamente divisível por $ (x - a) $ se, e somente se, \textbf{a} for uma de suas raízes. 

Sabendo disso, determine o valor real da constante \textbf{k} para que o polinômio geométrico $ V(x) = x^3 + kx^2 - 4x + 12 $ seja exatamente divisível pelo fator $ (x + 3) $.

\vspace{0cm} \noindent\textcolor{CorLinha}{\rule{\linewidth}{0.5pt}} \vspace{-0.2cm}
\subsubsection*{QUESTÃO 22}

Analise as afirmativas a seguir sobre as propriedades de grau nas operações de divisão de polinômios. Preencha com V (verdadeiro) ou F (falso) e justifique matematicamente as afirmativas consideradas falsas. Considere $ P(x) $ o dividendo e $ D(x) $ o divisor, com $ \text{grau}(P) \geq \text{grau}(D) $.

I) O grau do polinômio quociente será sempre igual à diferença exata entre o grau de $ P(x) $ e o grau de $ D(x) $.\\[0.3cm]
II) Na divisão de um polinômio de grau 5 por um polinômio de grau 2, o resto da divisão será obrigatoriamente um polinômio de grau 3.\\[0.3cm]
III) Se a divisão for exata, o polinômio resto é considerado identicamente nulo e, portanto, não possui grau definido.

\vspace{0cm} \noindent\textcolor{CorLinha}{\rule{\linewidth}{0.5pt}} \vspace{-0.2cm}
\subsubsection*{QUESTÃO 23}

(Estilo FUVEST) A teoria dos restos afirma que ao dividir um polinômio genérico $ P(x) $ por um divisor do segundo grau, o resto esperado será, no máximo, do primeiro grau, assumindo a forma $ R(x) = ax + b $. 

Um polinômio $ P(x) $ desconhecido deixa resto \textbf{4} quando dividido por $ (x - 1) $ e deixa resto \textbf{10} quando dividido por $ (x + 2) $. Utilizando essas informações e os princípios do Teorema do Resto, determine o polinômio resto da divisão de $ P(x) $ pelo produto $ (x - 1)(x + 2) $.

\vspace{0cm} \noindent\textcolor{CorLinha}{\rule{\linewidth}{0.5pt}} \vspace{-0.2cm}
\subsubsection*{QUESTÃO 24}

Para garantir o acoplamento perfeito em uma modelagem estrutural geométrica, o polinômio $ E(x) = x^4 - px^3 + qx^2 - 4x + 4 $ deve ser estritamente divisível pelo binômio composto $ x^2 - 4 $. 

Escreva o sistema linear gerado pelas condições do Teorema de D'Alembert e calcule os valores exatos dos coeficientes reais \textbf{p} e \textbf{q} que viabilizam esse acoplamento.

\vspace{0cm} \noindent\textcolor{CorLinha}{\rule{\linewidth}{0.5pt}} \vspace{-0.2cm}
\subsubsection*{QUESTÃO 25}

(ENEM - Adaptada) Uma indústria de embalagens projeta uma caixa retangular para o transporte de eletrônicos. A equipe de design estabeleceu que o volume total da embalagem é dado pelo polinômio $ V(x) = x^3 + 7x^2 + 14x + 8 $, e que a área da base da caixa é modelada por $ A(x) = x^2 + 5x + 4 $, onde \textbf{x} é uma medida de escala em centímetros.

\begin{center}
\begin{adjustbox}{max width=\linewidth}
\begin{tikzpicture}
% --- APLICANDO PALETAS MODERNAS E EFEITO 3D ---
\definecolor{caixafrente}{HTML}{E67E22}
\definecolor{caixatopo}{HTML}{F39C12}
\definecolor{caixalado}{HTML}{D35400}
\definecolor{textocor}{HTML}{2C3E50}

% LAYER 1: Sombra projetada
\fill[drop shadow={opacity=0.3, shadow xshift=3pt, shadow yshift=-3pt}, fill=white] (0,0) rectangle (0.1,0.1); % Hack para ativar sombra no escopo
\fill[black, opacity=0.15, rounded corners=2pt] (1.5, -0.2) -- (5.5, -0.2) -- (6.5, 1.3) -- (2.5, 1.3) -- cycle;

% LAYER 2: Face Lateral Direita
\filldraw[fill=caixalado, draw=white, thick, rounded corners=1pt] (4,0) -- (5.5,1.5) -- (5.5,3.5) -- (4,2) -- cycle;

% LAYER 3: Face Superior
\filldraw[fill=caixatopo, draw=white, thick, rounded corners=1pt] (0,2) -- (4,2) -- (5.5,3.5) -- (1.5,3.5) -- cycle;

% LAYER 4: Face Frontal
\filldraw[fill=caixafrente, draw=white, thick, rounded corners=1pt] (0,0) rectangle (4,2);

% LAYER 5: Textos e Cotas
\node[fill=white, inner sep=2pt, rounded corners=2pt, text=textocor, font=\bfseries\small] at (2,1) {$V(x) = x^3 + 7x^2 + 14x + 8$};
\node[fill=white, inner sep=1pt, rounded corners=2pt, text=textocor, font=\bfseries\footnotesize] at (2,2.7) {$A(x) = x^2 + 5x + 4$};

% Indicador de Altura
\draw[|<->|, >=stealth, thick, textocor] (-0.5,0) -- (-0.5,2) node[midway, fill=white, inner sep=1.5pt, font=\bfseries] {$H(x)$};
\end{tikzpicture}
\end{adjustbox}
\end{center}

Sabendo que o volume de um prisma retangular é o produto da área da base pela sua altura $ H(x) $, determine a expressão reduzida que representa a altura da embalagem efetuando a divisão polinomial. 

a) $ H(x) = x - 2 $\\[0.3cm]
b) $ H(x) = x + 1 $\\[0.3cm]
c) $ H(x) = x + 2 $\\[0.3cm]
d) $ H(x) = x + 3 $\\[0.3cm]
e) $ H(x) = x + 4 $

\vspace{0cm} \noindent\textcolor{CorLinha}{\rule{\linewidth}{0.5pt}} \vspace{-0.2cm}
\subsubsection*{QUESTÃO 26}

O Dispositivo de Briot-Ruffini pode ser aplicado consecutivamente para verificar a multiplicidade de uma raiz. Sabendo que o polinômio $ P(x) = x^4 - 5x^3 + 9x^2 - 7x + 2 $ admite a raiz $ x = 1 $, aplique o dispositivo iterativamente para provar se essa raiz é simples, dupla ou tripla. Demonstre os esquemas numéricos utilizados na sua verificação.

\vspace{0cm} \noindent\textcolor{CorLinha}{\rule{\linewidth}{0.5pt}} \vspace{-0.2cm}
\subsubsection*{QUESTÃO 27}

(Estilo UNICAMP) Na engenharia civil, o momento fletor ao longo de uma viga biapoiada sujeita a um carregamento irregular é descrito por um polinômio do terceiro grau $ M(x) $. Os apoios garantem que o momento fletor nas extremidades (pontos $ x = 0 $ e $ x = 6 $) seja nulo.

\begin{center}
\begin{adjustbox}{max width=\linewidth}
\begin{tikzpicture}
% --- APLICANDO PALETAS MODERNAS E SOMBREAMENTO ---
\definecolor{viga}{HTML}{95A5A6}
\definecolor{carga}{HTML}{3498DB}
\definecolor{apoio}{HTML}{7F8C8D}
\definecolor{linhaeixo}{HTML}{2C3E50}

% Carga distribuída (Setas)
\foreach \x in {0.5, 1.5, 2.5, 3.5, 4.5, 5.5} {
    \draw[->, >=stealth, thick, carga] (\x, 2) -- (\x, 0.7);
}
\draw[carga, thick] (0.2, 2) -- (5.8, 2);
\node[carga, font=\bfseries\small] at (3, 2.3) {Carregamento Irregular};

% Viga
\filldraw[fill=viga, draw=linhaeixo, thick, rounded corners=1pt, drop shadow={opacity=0.3}] (0,0.2) rectangle (6,0.6);

% Apoios (Triângulos)
\filldraw[fill=apoio, draw=linhaeixo, thick] (0.2, 0.2) -- (-0.2, -0.4) -- (0.6, -0.4) -- cycle;
\filldraw[fill=apoio, draw=linhaeixo, thick] (5.8, 0.2) -- (5.4, -0.4) -- (6.2, -0.4) -- cycle;

% Eixo x cotas
\draw[|<->|, >=stealth, thick, linhaeixo] (0.2, -0.8) -- (5.8, -0.8) node[midway, fill=white, inner sep=2pt, font=\bfseries] {$6\text{ m}$};
\node[font=\bfseries, linhaeixo] at (0.2, -1.2) {$x = 0$};
\node[font=\bfseries, linhaeixo] at (5.8, -1.2) {$x = 6$};

\end{tikzpicture}
\end{adjustbox}
\end{center}

Considerando que $ M(x) $ é divisível por $ x $ e também por $ (x - 6) $, e sabendo que no ponto $ x = 3 $ o momento medido foi de \textbf{54 kN.m}, determine a expressão algébrica completa de $ M(x) $ assumindo que o coeficiente do termo linear ($ x $) seja igual a zero.

\vspace{0cm} \noindent\textcolor{CorLinha}{\rule{\linewidth}{0.5pt}} \vspace{-0.2cm}
\subsubsection*{QUESTÃO 28}

Quando deparamos com expoentes massivos, a divisão em chaves é obsoleta. Considere o polinômio $ P(x) = x^{200} - 3x^{199} + x - 5 $. Calcule o resto da divisão desse polinômio gigantesco pelo binômio $ (x - 1) $. Justifique teoricamente qual o teorema que sustenta a sua resolução em apenas uma linha de cálculo.

\vspace{0cm} \noindent\textcolor{CorLinha}{\rule{\linewidth}{0.5pt}} \vspace{-0.2cm}
\subsubsection*{QUESTÃO 29}

(Estilo VUNESP) A divisão do polinômio $ P(x) = 2x^3 - mx^2 + nx - 3 $ por $ (x - 1) $ resulta em um resto igual a \textbf{2}. Sabe-se também que $ P(x) $ é divisível de forma exata por $ (x + 1) $. 

A partir do equacionamento dessas duas condições, estruture a matriz ou o sistema linear correspondente para encontrar os valores dos coeficientes \textbf{m} e \textbf{n}.

\vspace{0cm} \noindent\textcolor{CorLinha}{\rule{\linewidth}{0.5pt}} \vspace{-0.2cm}
\subsubsection*{QUESTÃO 30}

Sistemas de criptografia baseados em hardware utilizam a divisão de polinômios binários para validação de integridade (como o protocolo CRC). O esquema abaixo ilustra um bloco de processamento genérico onde uma mensagem polinomial $ M(x) $ sofre a ação de uma chave divisora $ C(x) $.

\begin{center}
\begin{adjustbox}{max width=\linewidth}
\begin{tikzpicture}
% --- APLICANDO PALETAS E SOMBRAS (Diagrama de Blocos) ---
\definecolor{blocomsg}{HTML}{3498DB}
\definecolor{blocochave}{HTML}{9B59B6}
\definecolor{blocoprocess}{HTML}{34495E}
\definecolor{blocoresto}{HTML}{E74C3C}
\definecolor{textobranco}{HTML}{FFFFFF}

% Mensagem M(x)
\node[fill=blocomsg, text=textobranco, font=\bfseries, minimum width=3cm, minimum height=1cm, rounded corners=3pt, drop shadow={opacity=0.2}] (msg) at (0,3) {Mensagem $M(x)$};

% Chave C(x)
\node[fill=blocochave, text=textobranco, font=\bfseries, minimum width=3cm, minimum height=1cm, rounded corners=3pt, drop shadow={opacity=0.2}] (chave) at (5,3) {Divisor $C(x)$};

% Processador (Divisão)
\node[fill=blocoprocess, text=textobranco, font=\bfseries, minimum width=4cm, minimum height=1.5cm, rounded corners=3pt, drop shadow={opacity=0.3}] (process) at (2.5,1) {Módulo de Divisão};

% Resto R(x) (Hash)
\node[fill=blocoresto, text=textobranco, font=\bfseries, minimum width=3cm, minimum height=1cm, rounded corners=3pt, drop shadow={opacity=0.2}] (resto) at (2.5,-1) {Resto Hash $R(x)$};

% Setas de fluxo
\draw[->, >=stealth, thick, black] (msg) -- (0, 1.5) -- (process.west);
\draw[->, >=stealth, thick, black] (chave) -- (5, 1.5) -- (process.east);
\draw[->, >=stealth, thick, black] (process) -- (resto) node[midway, fill=white, inner sep=2pt, font=\footnotesize\bfseries] {Extração};

\end{tikzpicture}
\end{adjustbox}
\end{center}

Seja a mensagem $ M(x) = x^5 - 2x^4 + 3x^3 - x^2 + 4 $ e a chave divisora $ C(x) = x^2 + 1 $. Determine o polinômio Hash $ R(x) $ gerado por esse bloco de processamento e explique a relevância de o grau de $ R(x) $ ser obrigatoriamente inferior ao grau de $ C(x) $.

\vspace{0cm} \noindent\textcolor{CorLinha}{\rule{\linewidth}{0.5pt}} \vspace{-0.2cm}
\subsubsection*{QUESTÃO 31}

Um polinômio $ P(x) $ do terceiro grau possui a forma $ P(x) = x^3 - 6x^2 + 11x - 6 $. Sem utilizar tentativa e erro livre, aplique o raciocínio dos divisores do termo independente (Teorema das Raízes Racionais) e o Dispositivo de Briot-Ruffini para encontrar o conjunto completo de raízes desse polinômio, detalhando as etapas do rebaixamento de grau.

\vspace{0cm} \noindent\textcolor{CorLinha}{\rule{\linewidth}{0.5pt}} \vspace{-0.2cm}
\subsubsection*{QUESTÃO 32}

Julgue as proposições a seguir, fundamentadas na Teoria dos Restos e na Divisibilidade Polinomial, preenchendo com V (verdadeiro) ou F (falso). Exija-se rigor absoluto na justificativa matemática para sentenças falsas.

I) Se um polinômio $ P(x) $ é divisível simultaneamente por $ (x - a) $ e por $ (x - b) $, com $ a \neq b $, então ele obrigatoriamente será divisível pelo produto $ (x - a)(x - b) $.\\[0.3cm]
II) O resto da divisão do polinômio $ x^n - a^n $ por $ (x - a) $ será nulo apenas se \textbf{n} for um número inteiro par.\\[0.3cm]
III) Ao aplicar o algoritmo da chave, se o divisor possui grau 3, o resto será rigorosamente um polinômio de grau 2 em todos os cenários possíveis.

\vspace{0cm} \noindent\textcolor{CorLinha}{\rule{\linewidth}{0.5pt}} \vspace{-0.2cm}
\subsubsection*{QUESTÃO 33}

(Estilo ITA) O limite do raciocínio analítico em polinômios envolve manipular suas raízes sem necessariamente conhecê-las inicialmente. Sabe-se que o polinômio $ P(x) = x^3 - 12x^2 + 44x - 48 $ possui três raízes reais distintas e que essas raízes formam uma Progressão Aritmética (P.A.). 

Utilizando as propriedades operacionais dos polinômios e a definição estrutural de uma P.A. de três termos $ (r - d, r, r + d) $, determine as três raízes do polinômio $ P(x) $.

\vspace{0cm} \noindent\textcolor{CorLinha}{\rule{\linewidth}{0.5pt}} \vspace{-0.2cm}
\subsubsection*{QUESTÃO 34}

Considere a expressão polinomial $ Q(x) = (x - 2)^5 + (x - 3)^4 - 1 $. Pela aplicação direta do Teorema do Resto, calcule qual será o resto quando $ Q(x) $ for dividido pelo binômio $ (x - 3) $. Mostre que a enorme potência algébrica não precisa ser expandida para solucionar o problema.

\vspace{0cm} \noindent\textcolor{CorLinha}{\rule{\linewidth}{0.5pt}} \vspace{-0.2cm}
\subsubsection*{QUESTÃO 35}

Um problema clássico de fatoração superior exige observar padrões. O polinômio $ S(x) = x^5 - 1 $ é dividido por $ (x - 1) $. O polinômio quociente $ Q(x) $ dessa divisão exata forma uma sequência geométrica em seus termos. 

Escreva a forma expandida completa do polinômio quociente $ Q(x) $ e, em seguida, calcule o valor numérico desse quociente quando avaliado em $ x = -1 $.

\end{multicols*}
\end{document}